\begin{abstract}

% Factual draft (2026-10-22), revised 2026-10-29: translation of the resumo,
% third person, no internal labels, de-identification rather than anonymization.
Classroom recordings expose students' and teachers' identities. This work
investigates whether replacing every detected face with a synthetic one reduces
automatic face re-identification while keeping behavioral signal (expression,
pose, emotion), measured on the video a viewer sees. The pipeline has three
stages: SCRFD-10GF detection with tracking and a recall post-pass; replacement
with BLANKET (one Stable Diffusion XL synthetic identity per track and an
\texttt{inswapper} face swap per frame); and fail-closed composition, which hides
every known face that has no valid replacement. The main contribution pushes the
swap's identity embedding away from the track's estimated real identity. On three
short videos (29\,s in total), held-out FaceNet rank-1 drops from
\resulten{e3c/pooled/native/facenet.rank1} to
\resulten{e3c/pooled/b-1.6/facenet.rank1} with no measurable expression or pose
cost; the direction of the effect holds when the push strength is chosen without
the evaluated video. With a privacy filter the final video reaches
\resulten{e3e/pooled/gate/facenet_final.rank1}, between
\resulten{derived/gate-final/x-chance-min} and
\resulten{derived/gate-final/x-chance-max} times chance depending on the video,
against \resulten{e2/pooled/blur/facenet_final.rank1} for blur, which erases
behavior. Gradient guidance inside the diffusion sampler and identity-code
optimization in a generative adversarial network have no measurable effect.
Running the main process on the GPU and adaptive frame sampling bring the
projected per-face loop cost to \resulten{e5/pooled/adapt-0.08/class-hours} GPU
hours per hour of class. No method keeps emotion near its noise floor. The
result is partial face de-identification, another point on the privacy--utility
trade-off, not anonymization: hair, head outline, clothing and voice remain.

\end{abstract}
